% !TeX root = ../main.tex

\chapter{基于注意力机制与图神经网络的交通流量预测}\label{chap:based_gnn_attention}
\section{引言}
尽管基于深度卷积神经网络的交通流预测算法在近年来取得了显著的成果, 尤其是图神经网络 (GNN) 的引入, 图神经网络在处理路网拓扑结构时具有明显优势, 能够有效捕捉节点间的复杂依赖关系, 极大增强了模拟交通数据空间特征的能力。 但是现有算法在对交通流量行为的非线性变换特征进行学习时仍显不足, 尤其是在捕捉数据时空特性方面。随着路网结构的不断扩展, 复杂的道路布局对模型提出了更高的要求。虽然图神经网络能够较好地适应这些拓扑结构, 并结合递归网络或时间卷积网络来有效学习时间维度的特征, 预测效果有所提高, 但仍有一些问题亟待解决。

\begin{enumerate}[(1)]
  \item 递归网络或时间卷积网络在处理长序列数据时常常遇到梯度消失或梯度爆炸的问题, 这会导致训练过程中的不稳定性。特别是在深层网络中, 这些问题变得尤为突出, 模型在时间维度上的学习能力有限, 从而影响了整体的预测效果。梯度消失或爆炸不仅使得网络难以有效更新权重, 而且可能导致信息在传递过程中逐渐丧失, 影响长期依赖关系的学习。

  \item 图神经网络虽然能够在空间维度上取得良好的预测效果, 但在交通流预测的过程中, 节点之间的流动性并非一成不变。在现有研究中, 通常通过固定的顺序或权重来定义节点之间的信息传递方式, 这种方法忽略了现实中路网节点间流量转移的概率差异。每个节点的流量转移能力可能受到多个因素的影响, 比如道路类型、 周围设施和天气等, 如何更加精准地表达这些节点间的异质性, 是未来研究的重要方向。
  
  \item 图卷积网络在处理交通流数据时, 可能会遭遇过平滑问题。过平滑是指随着图卷积层数的加深, 节点之间的信息过度融合, 导致模型难以区分不同节点的特征, 进而引发过拟合问题。特别是在深度图神经网络中, 过平滑会使得节点的表示变得过于相似, 从而影响模型的泛化能力, 导致预测效果不如预期。因此, 如何平衡图卷积层数与模型性能之间的关系, 避免过平滑的发生, 是当前图神经网络在交通流预测中的一个亟待解决的问题。
\end{enumerate}

针对上述三个问题, 论文使用基于时空注意力机制和图网络模型的交通流预测方法。算法基于图神经网络的交通流预测模型, 设计了一种基于时间与空间注意力增强的图卷积交通流预测模型, 序列的学习将不再使用循环神经网络, 改为同时使用时间维度上的注意力机制与卷积操作, 同时为了区分邻域节点对于中心节点的贡献程度, 采用了图注意力机制对于连接节点的边赋予不同的权重, 提高了预测任务的精度。下文详细介绍本章模型的结构、原理和实验结果。


\section{基于时空注意力机制与图网络结合的交通流量预测方法设计思路}

算法旨在从交通流数据中更好地提取出数据的时空特征, 整个模型如图 \ref{fig:model_structure} 所示。模型主要分为两个部分: (1)动态时空相关性信息捕获，通过时空注意机制得以实现。复杂空间的相关性信息由空间注意获取。动态时间相关性，则依赖于时间注意力机制在不同时间片段上的运用。邻接矩阵的信息与空间注意力机制，在图卷积层的卷积运算中融合。(2)时空卷积, 使用图卷积来捕获空间模式, 并使用时间卷积来描述时间特征。另外, 考虑到交通流量数据在每天的不同时间段, 每周的不同时间段都具有不同的变化特征, 故选取样本片段时, 将选取预测时间段最近片段、昨日同期片段和上周同期片段, 并对三个分量的输出进行加权融合, 得到动态预测结果。


  
\begin{figure}
  \begin{center}
    \includegraphics[width=0.7\textwidth]{01960c08-a17d-7162-9acb-7f9eeb4038d6_17_420_1050_854_761_0.jpg}
  \end{center}
  \caption{模型结构示意图 \parencite{2019_Guo_Attention_based_spatial_temporal_graph_convolutional_networks_for_traffic_flow_forecasting}}
  \label{fig:model_structure}
\end{figure}

针对上文中提到的问题, 模型在频域卷积的过程中加入了注意力机制, 具体添加方式为: 在每次图卷积之前加入图注意力层, 通过学习节点之间的相似性计算得到相似性矩阵, 再将节点相似性矩阵与图卷积核进行哈德玛乘积, 得到带有注意力信息的图卷积核运算矩阵。 再通过时间卷积学习时间特征, 最后通过全连接层, 输出所需预测的时间片段。在本文所主要使用的时空注意力机制-图卷积模型中, 主要即堆叠 2 个时空块和 1 个全连接层组成。在每个时空块中都有一个时空注意模块和一个时空卷积模块。为了优化训练效率, 本文中也采用了残差学习框架 \parencite{2016_He_Deep_residual_learning_for_image_recognition}。


\section{模型具体结构分析}

\subsection{时空注意力机制}

交通流量预测问题本质上属于时间序列预测范畴, 需要依赖历史时段数据对未来状态进行推断。特殊之处在于, 研究对象具有显著区别于常规时间序列的特征。若仅采用序列到序列的预测方式, 道路网络内在拓扑结构信息将难以被有效捕获。由此可见, 采用图网络模型与空间注意力机制相结合的方法更为适宜, 能够实现对道路空间关系的建模。现有研究中关于时间维度预测的处理方式, 多数情况下可见循环神经网络或时序卷积网络的应用实例。本文使用的时间注意力机制及时间卷积可以解决上述循环时间网络所出现的一些问题。

在聚合邻居节点的过程中，基于注意力机制，图注意力网络被采用。由于针对不同邻居节点的权重差异性，图神经网络模型的表达能力因此获得显著提升。权重的动态调整过程呈现出明显的空间异质性格局。由此可见，注意力机制的引入使得节点间的信息传递效率产生了质的飞跃，特别是在处理异构图的场景下表现得尤为突出。注意力机制最初受到人类视觉聚焦方式的启发, 能够让模型在处理信息时自动聚焦于重要部分, 从而提高效率和准确性。类似于人类在视觉感知中会聚焦于某个区域, 忽略其余部分, 注意力机制通过为输入的各个部分分配不同的权重, 使模型能够在复杂任务中自动关注关键信息。这一机制最早在自然语言处理领域中得到了广泛应用, 尤其是在机器翻译和文本生成等任务中, 帮助模型在不同的词语和句子之间建立起关联, 从而改善理解与生成的效果。

在深度学习中, 注意力机制通过计算输入信息的加权和, 使得模型能够在特定的上下文中获取更多相关信息, 特别是在长序列的任务中展现出显著优势。它不仅提高了模型的表达能力, 也大大提升了计算效率, 成为了图像处理、语音识别等多个领域的重要组成部分。\textcite{2015_Bahdanau_Neural_machine_translation_by_jointly_learning_to_align_and_translate} 也首次在 Seq2Seq 模型中提出注意力机制, 注意力机制旨在解决 RNN 循环结构在处理长序列时会出现的信息丢失问题, 通过优化序列之间的交互方式, 增强了模型对长距离依赖关系的捕获能力。

注意力机制的数学表述如下:

在注意力机制中，Source（源）通常指的是输入序列或特征。这些是模型根据其重要性进行权重分配的对象。源可以是一个单词、特征向量或任何其他类型的输入数据。Query（查询）是模型在注意力机制中用来查找相关信息的部分。它通常是当前需要进行处理或预测的目标。比如，在机器翻译中，查询可能是当前要翻译的单词的上下文。在注意力机制中，每个源通常会与一个Key和一个Value相关联。Key用于与查询进行匹配或计算相似度，而Value则包含了源的实际信息或特征。最终的注意力输出是根据Key的匹配度加权求和的Value。Attention Value（注意力值）是通过将Query与各个Key进行比较（如计算相似度或匹配度），得到的每个Source的相关性分数。这些分数表明了模型在当前预测中应关注的源的程度。高分表示源对当前查询的贡献较大，低分表示贡献较小。注意力机制的定义如下:

Attention(Query, Source) \(= \sum\) similarity(Query, Key) \(\cdot\) Value.

在图注意力网络（Graph Attention Network, GAT）中，Source、Query、Attention Value 和 Key-Value 的作用与标准注意力机制相似，但它们应用于图结构数据。具体来说，Source 是图中每个节点的特征向量，表示节点的输入信息；Query 是当前节点的查询向量，用于与邻接节点的特征进行匹配，来决定哪些节点对当前节点的表示贡献较大；Key 是每个邻居节点的特征向量，用于与查询向量进行匹配，计算相似度；而 Value 是节点本身的特征信息，用于在计算加权平均时作为输出。因此，图注意力机制通过对每个节点的邻居节点进行加权聚合，来学习节点之间的关系和图的结构信息。

假设中心节点为 \({{v}}_{{i}}\) , 我们设邻居节点 \({{v}}_{{j}}\) 到 \({{v}}_{{i}}\) 在某一层后的权重系数为:
\begin{equation}
{e}_{i, j} = a\left( {W{h}_{i}, W{h}_{j}}\right).
\end{equation}

\({W} \in  \mathbb{R}^{d\left( {l + 1}\right)  * d\left( l\right) }\) 是对该层节点特征进行仿射变换的权重参数。$a()$是计算两个节点相关度的函数。关于 \({a}\) 的选择, 可以灵活多样，输出节点之间的相似度即可。既可以选择对进行仿射变换后的节点特征进行内积运算，也可以通过网络层的再次嵌套。本文选择的方法为后者：
\begin{equation}
{{e}}_{{i}, {j}} = \text{ Leaky }\operatorname{ReLU}\left( {{{a}}^{{T}}\left\lbrack  {{{{Wh}}}_{{i}}\parallel {{{Wh}}}_{{j}}}\right\rbrack  }\right).
\end{equation}

其权重参数 \({a} \in  {\mathbb{R}}^{2 * {d}\left( {{l} + 1}\right) }\) , 激活函数设计为 LeakyReLU (   )。为了更好地分配权重, 需要将与所有邻居计算出的相关度进行统一的归一化处理, 具体形式为 softmax 归一化:
\begin{equation}
{\alpha }_{i, j} = {\operatorname{softmax}}_{j}\left( {e}_{i, j}\right)  = \frac{\exp \left( {e}_{i, j}\right) }{\mathop{\sum }\limits_{{{v}_{k} \in  N\left( {v}_{i}\right) }}\exp \left( {e}_{ij}\right) }.
\end{equation}

一旦完成上述权重系数的计算, 按照注意力机制加权求和的思路, 节点 \({{v}}_{{i}}\) 新的特征向量为:
\begin{equation}
{h}_{i}^{\prime } = \sigma \left( {\mathop{\sum }\limits_{{{v}_{j} \in  N\left( {v}_{i}\right) }}{\alpha }_{i, j}W{h}_{j}}\right).
\end{equation}

关于空间的注意力机制计算如下:
\begin{equation}
S = {V}_{s} \cdot  \sigma \left( {\left( {{\mathcal{X}}_{h}^{\left( r - 1\right) }{W}_{1}}\right) {W}_{2}{\left( {W}_{3}{\chi }_{h}^{\left( r - 1\right) }\right) }^{T} + {b}_{s}}\right)
\end{equation}
\begin{equation}
{S}_{i, j}^{\prime } = \frac{\exp \left( {S}_{i, j}\right) }{\mathop{\sum }\limits_{{j = 1}}^{N}\exp \left( {S}_{i, j}\right) }.
\end{equation}

关于时间的注意力机制计算如下:
\begin{equation}
{E} = {{V}}_{{e}} \cdot  \sigma \left( {\left( {{\mathcal{X}}_{{h}}^{\left( {r} - 1\right) }{{U}}_{1}}\right) {{U}}_{2}{\left( {{U}}_{3}{\mathcal{X}}_{{h}}^{\left( {r} - 1\right) }\right) }^{{T}} + {{b}}_{{e}}}\right)
\end{equation}
\begin{equation}
{E}_{i, j}^{\prime } = \frac{\exp \left( {E}_{i, j}\right) }{\mathop{\sum }\limits_{{j = 1}}^{T}\exp \left( {E}_{i, j}\right) }, 
\end{equation}
其中 \({{V}}_{{s}}\) 、 \({{W}}_{1}\) 、 \({{W}}_{2}\) 、 \({{W}}_{3}\) 、 \({{b}}_{{s}}\) 、 \({{V}}_{{e}}\) 、 \({{U}}_{1}\) 、 \({{U}}_{2}\) 、 \({{U}}_{3}\) 、 \({{b}}_{{e}}\) 均为可学习的参数.

时空注意模块让网络自动地相对较多地关注有价值的信息。从计算的层面来说, 将由时间注意机制调整的输出作为输入, 输入到空间卷积模块中。具体来说, 为了动态调整节点之间的相关性, 对于图卷积运算, 将图卷积参数矩阵 \({M}\) 融合空间注意矩阵 \({S}_{i, j}\) 信息, 然后得到 \({M} \odot  {S}_{i, j}^{\prime }\) . 时空卷积模块主要用来建模交通数据的时空依赖性。时空卷积模块包括用于从原始的基于图的交通网络结构中捕获空间特征的图卷积, 以及用于描述来自附近的时间切片的依赖关系的时间维度上的时间卷积。在空间维度上的图卷积操作捕获了每个节点的相邻信息后, 在时间维度上进一步叠加标准卷积层, 通过合并相邻时间片上的信息来更新节点的信号, 最后通过全连接层连接输出预测信号。


\subsection{图滤波器优化}


最后, 提出关于图卷积层中图滤波器的优化. 图谱滤波的思想是调制图信号的频率, 使得其中一些频率分量被保留或放大, 而另一些频率分量被移除或减小。 因此, 给定一个图信号, 首先需要对其进行图傅里叶变换, 以获得它的图傅里叶系数, 然后对这些系数进行调制, 再在图域中重构该信号。

根据第二章图基础理论介绍, 图傅里叶变换被定义为:
\begin{equation}
\widehat{f} = {U}^{T}f, 
\end{equation}
由信号 \(f\) 所进行的图傅里叶变换操作中，其生成系数可由 \(\widehat{f}\) 表示。对于图信号 \(f\) 而言，各个图傅里叶分量所具有的贡献度，正是通过这些傅里叶系数得以体现。对应于第 \({i}\) 个图傅里叶分量 \({u}_{i}\) 的，是 \(\widehat{f}\) 当中第 \({i}\) 个元素的特定存在。实例表明，当需要对图信号 \(f\) 的频率特性实施调制时，可采取对图傅里叶系数执行滤波处理的方式。以频率 \({\lambda }_{i}\) 作为输入变量的函数 \(\gamma \left( {\lambda }_{i}\right)\) ，其具体形式决定了相应频率分量将被如何调制。
\begin{equation}
\widehat{f} = \widehat{f}\left\lbrack  i\right\rbrack   * \gamma \left( {\lambda }_{i}\right) , i = 1, 2\ldots N.
\end{equation}
矩阵形式表示如下:
\begin{equation}
{\widehat{f}}^{\prime } = \gamma \left( \Lambda \right)  * \widehat{f} = \gamma \left( \Lambda \right) {U}^{T}f.
\end{equation}
有了滤波系数之后, 就可以将图信号重新构建到图域:
\begin{equation}
{f}^{\prime } = U{\widehat{f}}^{\prime } = U * \gamma \left( \Lambda \right) {U}^{T}f.
\end{equation}

滤波过程可视为将算子 \(U * \gamma \left( \Lambda \right) {U}^{T}\) 应用于输入的图信号。 \(\gamma \left( \Lambda \right)\) 控制着图信号 \(f\) 的每个频率分量的滤波过程。例如, 如果 \(\gamma \left( {\lambda }_{i}\right)\) 等于 0, 即 \(\widehat{f}\left\lbrack  {i}\right\rbrack   = 0\) , 则频率分量 \({u}_{i}\) 会从图信号中移除。如果知道该如何调制输入信号中的频率, 那么就可以用相应的方式设计函数 \(\gamma \left( \Lambda \right)\)。然而, 在实际使用中, 当把基于谱的滤波器作为图滤波器时, 哪些频率更重要往往是未知的。 因此, 就像传统的神经网络一样, 可以通过数据驱动的方式学习合适的图滤波器。具体来说, 可以用某些函数对 \(\gamma \left( \Lambda \right)\) 建模, 然后在数据的监督下学习参数。一种自然的尝试是在设计 \(\gamma \left( \cdot \right)\) 时给予它充分的自由。 具体而言, 函数 \(\gamma \left( \cdot \right)\) 定义为:
\begin{equation}
\gamma \left( {\lambda }_{i}\right)  = \left\lbrack  \begin{matrix} {\theta }_{1} & \cdots & 0 \\  \vdots &  \ddots  & \vdots \\  0 & \cdots & {\theta }_{N} \end{matrix}\right\rbrack.
\end{equation}

此类非参数化卷积算子具有显著的局限性。节点数\({N}\)与需学习参数数量相等，当\({N}\)取值较大时，内存占用量将显著提升，数据拟合需求亦随之增加。稠密矩阵特性可能存在于卷积算子\(U * \gamma \left( \Lambda \right) {U}^{T}\)中，因此输出信号\({f}^{\prime }\)的第\({i}\)个元素计算过程可能涉及图中全部节点。拉普拉斯矩阵特征分解的包含于该算子计算过程中，加之稠密矩阵相乘运算的存在，计算代价高昂。

Poly-filter \parencite{2021_Han_Dynamic_and_multi_faceted_spatio_temporal_deep_learning_for_traffic_speed_forecasting} 旨在解决谱滤波问题。谱滤波在 GNN 中至关重要, 它使得模型能够有效地处理图结构数据, 从而更好地学习节点和边的表示。Poly-filter 特别设计用于通过应用多项式滤波器来平滑图信号的谱响应。这一方法建立在谱卷积框架的基础上, 其中图拉普拉斯算子在定义图结构中起着核心作用。通过使用多项式滤波器, Poly-filter 能够更高效地处理图数据, 避免了通常在谱方法中需要进行昂贵的特征值分解。

然而, Poly-filter 也存在一些局限性。其主要挑战之一是滤波器的频率响应在大图或者多项式度数较高时可能会不稳定。多项式的度数直接影响滤波器捕捉图中长程依赖的能力, 但也增加了过度平滑的风险, 即模型可能失去准确预测所需的细粒度信息。此外, 主要问题在于多项式基（如 \(1, {x}, {{x}}^{2}, \ldots\)），这些系数间存在相互依赖性，学习过程中的扰动易引发不稳定性。一个系数的更新行为可能导致其余系数发生改变，由此可见基函数正交化的重要性。

为了解决这些问题, \textcite{2016_Defferrard_Convolutional_neural_networks_on_graphs_with_fast_localized_spectral_filtering} 提出了 Cheby-Filter 作为替代方案。 Cheby-Filter 提供了一种减少过度平滑和计算复杂度的解决方案。它们利用 Chebyshev 多项式, 能够高效地近似图卷积, 而无需显式地计算图拉普拉斯算子的特征值。Cheby-Filter 还提供了更好的控制, 能够稳定地处理图的不同大小。由于它们能够以较少的计算资源处理长程依赖, 因此在大规模图学习任务中, Cheby-Filter 被认为是 Poly-filter 的一个有前景的替代方案。

切比雪夫多项式可以由如下递推关系生成:
\begin{equation}
{T}_{k}\left( y\right)  = {2y}{T}_{k - 1}\left( y\right)  - {T}_{k - 2}\left( y\right), 
\end{equation}
其中 \({y\epsilon }\left\lbrack  {-1, 1}\right\rbrack  , {T}_{0}\left( y\right)  = 1, {T}_{1}\left( y\right)  = y\)。切比雪夫多项式彼此正交, 由于切比雪夫多项式的定义域是 \(\left\lbrack  {-1, 1}\right\rbrack\) , 为了利用切比雪夫多项式近似卷积算子, 需要重新缩放并平移拉普拉斯矩阵的特征值, 具体如下:
\begin{equation}
\widetilde{{\lambda }_{l}} = \frac{2 * {\lambda }_{l}}{{\lambda }_{\max }} - 1.
\end{equation}
因此，$y$为矩阵时即为：
\begin{equation}\tilde{\boldsymbol{L}}=\frac{2\boldsymbol{L}}{\lambda_{\max}}-\boldsymbol{I}.\end{equation}


\section{实验结果分析}
本节对于上文提出的基于时空注意力机制-图网络交通流预测模型进行实验分析验证。实验在多个数据集上展开, 包括 METR-LA 和 SZ 数据集, 下文将详细介绍。实验主要包含：由模型和其他主流预测模型的横向对比结果, 以及模型各部分模块的消融实验对比。

\subsection{实验细节设置}


本文基于 pytorch 框架实现了时空注意力机制-图卷积模型。\textcite{2016_Kipf_Semi_supervised_classification_with_graph_convolutionalnetworks} 等检验了切比雪夫多项式 \({K} \in  \{ 1, 2, 3\}\) 的项的个数。但随着 \({K}\) 的增大, 预测性能提高的效率微乎其微。考虑到计算效率和预测性能的提高程度, 将 \({K}\) 设置为 3 以及沿时间维度的卷积核大小设置为 3。在模型训练之前, 需要对原始数据集进行预处理。根据传感器的地址信息与连通属性构建路网的图结构表示形式, 即构建邻接矩阵。为了捕捉交通流量的近期特征、日周期特征以及周周期特征, 这里选择了一个包含 \({36}/{12}\) 个时间戳的时间窗口作为预测输入。(LA 数据集为 36 个时间点, SZ 数据集为 12 个时间点)具体来说, 输入包括: 前一个小时的流量数据、前一日相同时间段的流量数据, 以及前一周相同时间段的流量数据。通过这种方式, 模型能够学习到交通流量在不同时间尺度下的变化模式。数据输出为预测未来一小时的交通节点的流量(LA 数据集为 12 个时间点, SZ 数据集为 4 个时间点)。在本章的实验中, 每个 Batch 设置大小为 32, 共 100 个 epoch。网络模块中时空注意力机制的相关映射维度为 64 维。所有的时间卷积层都使用 64 个卷积核。为方便对齐, 图卷积的输出维度也设置为 64 维。训练过程中, 初始学习率设置为 0.01, 使用 Adam 进行优化器进行训练, 每个 epoch 都对训练模型进行验证。将选择验证集损失最小时对应的模型参数, 作为最终模型参数, 用于在测试集上进行预测。按照 $6:2:2$ 的比例划分为训练集、验证集和测试集。

本文将主要围绕下列 2 个数据集展开实验:

\begin{enumerate}[(1)]
  \item METR\_LA (Metro Los Angeles) : 数据集是一个由加利福尼亚大学洛杉矶分校 (UCLA) 发布的交通流量数据集, 专门用于时间序列预测和交通流量预测任务。这个数据集为研究交通预测、时间序列分析、深度学习方法等领域提供了一个非常重要的基准。METR\_LA 数据集中的数据来源于洛杉矶地区的高速公路网络, 记录了在不同地点的交通流量信息。具体来说, 它包含了从 207 个传感器收集到的从 2012/03/01-2012/06/27 (约 34272 条记录) 的数据。这些传感器监测了汽车的流量、速度、占用率等信息, 且数据以时间序列的方式进行记录。这里使用行驶速度作为预测的变量, 行驶速度基本能反映交通拥堵情况.

  \item SZ 数据集: SZ-taxi 数据集是从深圳市的出租车运行数据系统采集而来, 在 T-GCN 网络中首次使用。数据集主要采集地区在罗湖区内, 数据类型是出租车的行驶速度。数据集的时间跨度为 2015 年 1 月 1 日到 2015 年 2 月 4 日, 相邻两个采样之间间隔 15 分钟。数据集内部包括了 156 个道路节点在不同时刻的速度值, 以及代表节点之间连接情况的邻接矩阵。
\end{enumerate}

另外, 这两个数据集也有不同的特点: SZ 数据集是在城市道路高度复杂化的区域采集而成, 其受到的外部因素的影响程度高于LA 高速公路数据集。


\subsection{实验结果}


预测结果选择 RMSE(均方根误差)、MAE(平均绝对误差)、和 MAPE(平均绝对百分比误差)作为最终的评估指标。最终各模型预测结果见下表。

RMSE 是预测值与实际值之间误差的平方的均值的平方根, 公式为:
\begin{equation}
\text{ RMSE } = \sqrt{\frac{1}{{n}}\mathop{\sum }\limits_{{{i} = 1}}^{{n}}{\left( {{y}}_{{i}} - {\widehat{{y}}}_{{i}}\right) }^{2}}.
\end{equation}

RMSE 对较大误差赋予更高的权重, 因为误差被平方后放大了, 因此它在处理具有较大异常值的情况时更加敏感。它通常用于强调对大误差的惩罚, 适用于对极端误差非常敏感的情形。

MAE 是预测值与实际值之间绝对误差的平均值, 公式为:
\begin{equation}
\text{ MAE } = \frac{1}{{n}}\mathop{\sum }\limits_{{{i} = 1}}^{{n}}\left| {{{y}}_{{i}} - {\widehat{{y}}}_{{i}}}\right|.
\end{equation}
MAE 反映了预测值偏离实际值的程度, 越小表示模型越准确。它不考虑误差的方向, 仅关心误差的大小, 因此适用于需要量化误差大小的情境。

MAPE 是误差的绝对值与实际值之比的平均值, 公式为:
\begin{equation}
\text{ MAPE } = \frac{1}{{n}}\mathop{\sum }\limits_{{{i} = 1}}^{{n}}\left| \frac{{{y}}_{{i}} - {\widehat{{y}}}_{{i}}}{{{y}}_{{i}}}\right|.
\end{equation}
MAPE 提供了一个百分比误差, 能够衡量预测的相对精度。由于它基于实际值与预测值的比例, 通常用于评估预测值在不同数量级的实际数据下的相对误差。然而, 如果实际值接近零, MAPE 会表现出较大的误差, 这在某些情形下可能导致不准确的评估。

为了验证本文提出的时空注意机制的优势, 使用了一些基线模型, 长短记忆神经网络、图嵌入技术、TGCN 模型等进行对比, 以展示其先进性。还设计了相关退化版本的模型, 分别为时空注意力模型（相关参数相同, 但不使用图卷积来学习节点间的相关性）, 时间注意力-图卷积模型（不使用图注意力机制）, 以验证图注意力机制能缓解图卷积过平滑的问题. 

由实验结果（见表 \ref{tab:SZ_results} 和表 \ref{tab:LA_results}）可知, 本文主要使用的时空注意力机制-图卷积模型（下列简称ASTGCN模型）, 预测精度较基线模型都有一定程度的提升。ASTGCN在LA数据集上MAE和RMSE分别较对比模型下降至少40.70\%和33.08\%. 在SZ数据集上MAE和RMSE分别较对比模型下降至少17.22\%和11.76\%.
\begin{table}[htbp]
  \caption{SZ 数据集实验结果}
  \label{tab:SZ_results}
  \centering
  \begin{tblr}{
      hline{1, Z} = {1pt}, 
      hline{2}, 
      width = \textwidth, 
      colspec = {X[2, c]X[1, c]X[1, c]X[1, c]}, 
    }
    SZ 数据集 & RMSE & MAE & MAPE \\
    T-GCN & 4.59 & 3.31 & 42.97\% \\
    AST-GCN & 4.05 & 2.74 & 9.21\% \\
    Graph\_Embedding & 23.47 & 13.69 & 13.50\% \\
    LSTM & 11.84 & 5.31 & 21.92\%
  \end{tblr}
\end{table}
\begin{table}[htbp]
  \caption{LA 数据集实验结果}
  \label{tab:LA_results}
  \centering
  \begin{tblr}{
      hline{1, Z} = {1pt}, 
      hline{2}, 
      width = \textwidth, 
      colspec = {X[2, c]X[1, c]X[1, c]X[1, c]}, 
    }
    LA数据集 & RMSE & MAE & MAPE \\
    T-GCN & 12.57 & 8.3 & 290.61\% \\
    AST-GCN & 6.94 & 3.57 & 10.22\% \\
    Graph\_Embedding & 31.49 & 17.28 & 47.00\% \\
    LSTM & 10.37 & 6.02 & 301.53\% \\
  \end{tblr}
\end{table}

使用传统的 LSTM 模型时, MAPE 较高, 说明该模型未能充分捕捉到即将发生拥堵的路口信息。LSTM 主要关注时间序列数据的时间依赖性, 但在处理交通流量预测时, 缺乏对空间信息和路网中节点间交互的有效建模。交通流量不仅是时间序列的动态变化, 还受到路网结构和相邻路口流量的影响, 因此 LSTM 的空间建模能力较弱, 无法及时捕捉到不同路口之间的复杂依赖关系, 从而影响了对拥堵情况的准确预测。

图嵌入模型在预测任务中的表现较差, 可能的原因在于它们通过预先学习的低维嵌入来表示节点或图结构, 这些嵌入虽然能够捕捉一定的结构信息, 但通常无法直接针对具体的下游任务进行优化。因此, 图嵌入在转化为具体任务时, 其学习到的表示可能对预测并不起作用, 缺乏足够的任务相关性。而与之相比, 图卷积网络采用端对端的学习方式, 能够在训练过程中自动调整图结构的表示, 从而更好地适应任务需求。GCN 通过消息传递机制动态更新节点特征, 并能够逐层学习更加精细的局部和全局图信息, 因此其在捕捉节点间复杂关系和结构依赖方面更为有效, 通常能取得更好的预测性能。

使用 T-GCN 预测模型的指标均不及本章节中的时刻注意力机制-图网络模型, 二者模型主要区别在, 前者使用循环神经网络学习序列特征, 后者使用时间注意力机制联合时间卷积学习序列特征。T-GCN 在 SZ 数据集上的实际值与预测值比较, T-GCN 对交通流量的升降趋势预测有一定的延迟, T-GCN 模型无法即使预测出即将发生堵塞的节点。图 \ref{fig:sz-test} 展示了 T-GCN 与本文使用的 ASTGCN 模型在 SZ 数据集测试集部分的预测效果, 本文的模型基本能够拟合真实交通流量, 而 T-GCN 在验证集预测上几乎失效。

\begin{figure}[htbp]
  \centering
  \includegraphics[width=0.9\textwidth]{01960c08-a17d-7162-9acb-7f9eeb4038d6_26_307_1179_1036_497_0.jpg}
  \includegraphics[width=0.9\textwidth]{01960c08-a17d-7162-9acb-7f9eeb4038d6_27_248_335_1157_553_0.jpg}
  \caption{T - GCN 与 ASTGCN 在 SZ 测试集上预测效果对比}
  \label{fig:sz-test}
\end{figure}
% \begin{figure}[htbp]
%   \centering
%   % 第一个子图
%   \begin{subfigure}[b]{\textwidth}
%       \centering
%       \includegraphics[width=0.8\textwidth]{01960c08-a17d-7162-9acb-7f9eeb4038d6_26_307_1179_1036_497_0.jpg}
%       \caption{T - GCN 与 ASTGCN 在 SZ 测试集上预测效果对比}
%       \label{fig:sub1}
%   \end{subfigure}
%   % 第二个子图
%   \begin{subfigure}[b]{\textwidth}
%       \centering
%       \includegraphics[width=0.9\textwidth]{01960c08-a17d-7162-9acb-7f9eeb4038d6_27_248_335_1157_553_0.jpg}
%       \caption{T - GCN 与 ASTGCN 在 SZ 测试集上预测效果对比}
%       \label{fig:sub2}
%   \end{subfigure}
% \end{figure}


图 \ref{fig:la-test} 展示了 T-GCN 与 ASTGCN 在 LA 数据集的预测精度对比, 可见 ASTGCN 在 LA 数据集上表现比 SZ 数据集更好, 更能预测出那些即将发生拥堵的交通节点。 因此, ASTGCN 模型更适用于受外部因素干预较小的交通网络数据。
\begin{figure}[htbp]
  \centering
  \includegraphics[width=1.0\textwidth]{01960c08-a17d-7162-9acb-7f9eeb4038d6_27_278_1220_1189_480_0.jpg}
  \caption{T - GCN 与 ASTGCN 在 LA 测试集上预测效果对比}
  \label{fig:la-test}
\end{figure}

消融实验相关结果（见表 \ref{tab:SZ_ablation_study_results} 和表 \ref{tab:LA_ablation_study_results}）证明，模型中一些模块的优化能够提升模型精度。

\begin{table}[htbp]
  \caption{SZ 数据集消融实验结果}
  \label{tab:SZ_ablation_study_results}
  \centering
  \begin{tblr}{
      hline{1, Z} = {1pt}, 
      hline{2}, 
      width = \textwidth, 
      colspec = {X[2, c]X[1, c]X[1, c]X[1, c]}, 
    }
    SZ 数据集 & RMSE & MAE & MAPE \\
    AST-GCN & 4.05 & 2.74 & 9.21\% \\
    w/o SpatialAttention & 4.07 & 2.80 & 8.46\% \\
    w/o GCN & 3.99 & 2.77 & 7.65\% \\
  \end{tblr}
\end{table}
\begin{table}[htbp]
  \caption{LA 数据集消融实验结果}
  \label{tab:LA_ablation_study_results}
  \centering
  \begin{tblr}{
      hline{1, Z} = {1pt}, 
      hline{2}, 
      width = \textwidth, 
      colspec = {X[2, c]X[1, c]X[1, c]X[1, c]}, 
    }
    LA数据集 & RMSE & MAE & MAPE \\
    AST-GCN & 6.94 & 3.57 & 10.22\% \\
    w/o SpatialAttention & 7.37 & 4.03 & 11.34\% \\
    w/o GCN & 6.87 & 3.82 & 12.86\% \\
  \end{tblr}
\end{table}

ASTGCN 模型比仅使用时间注意力机制的模型在 LA 数据集表现最好, MAE下降 11.41\%, RMSE 下降 5.83\%。这一提升主要归因于空间注意力机制的引入, 有效缓解了图卷积网络中的过平滑问题。空间注意力机制能够动态调整节点间的相互关系, 使得信息传播更加精细化, 避免了信息过度平滑, 从而改善了模型的表达能力和预测精度, 尤其在处理图数据时, 能够捕捉到更加局部且重要的特征。

在本章节中, 我们展示的算法在部分数据集上的结果表现逊色于消融 GCN 的退化模型。这一现象的背后原因, 可能与 GCN 中预先定义的图结构不完全适用于特定问题相关。传统的 GCN 方法通常依赖于固定的图结构, 该结构是在训练开始前就已经确定。然而, 许多实际问题中, 图结构并非固定, 而是具有一定的动态性或者应根据数据本身的特征来进行调整。例如, GWnet \parencite{2019_Wu_Graph_Wavelet_Networks_for_Deep_Learning_on_Graphs} 就采用了非固定结构的图网络, 这使得其能够在数据集的图结构变化时灵活地进行调整, 从而可能获得更好的性能表现。本文将在后续部分, 探索如何改进 GCN 方法, 通过引入更加灵活和动态的图结构, 以提升模型在不同任务中的适应性和表现。








\section*{3.4.3 多信息融合适应性}

现有的交通流研究主要依赖于数据集中的传感器数据, 预测模型通过历史交通数据来预测未来的流量。然而, 交通流量数据具有非线性特征, 导致其未来状态不仅仅与历史数据相关, 还会受到其他外部因素的影响, 比如道路类型、周边设施的种类以及天气条件等 \parencite{2019_Zhang_Flow_prediction_in_spatio_temporal_networks_based_on_multitask_deep_learning}。因此, 如何将这些外部信息整合到预测模型中, 并以合适的方式将其加入模型, 是一个值得进一步探讨的课题。因此, 本部分内容将论述时空注意力机制-图网络模型对于融合多信息特征的适应性。

外部信息源主要有两部分, 静态信息例如 POI (道路兴趣点指数), 通常指的是在特定区域内具有某种吸引力或重要性的地理位置。这些地点可能是旅游景点、 商店、餐馆、医院、学校等。POI 数据广泛应用于地图服务、地理信息系统(GIS)、 推荐系统等领域, 可以帮助分析和理解人们的活动模式、行为趋势等。POI 信息与交通流数据结合在一起, 可以为交通流预测、城市规划、推荐系统等应用提供有价值的辅助数据。通过分析 POI 数据与交通流之间的关系, 研究人员能够更好地理解交通模式以及外部因素对交通流的影响。动态信息内容随着时间的改变而改变。在多个动态信息中, 本文选取了天气信息作为外源特征。其他还包括道路占有率、车辆数等。

现有研究融合多信息的思路为直接将外源信息于原始信息拼接 \parencite{2023_胡晓荣_基于图神经网络的路网短时交通流预测算法研究}。而在本文中, 网络层的主要设计思路为, 在通过时空注意力层、时空卷积的所有运算后, 融合外部信息, 将外部信息进行仿射变换与交通序列的维度进行对齐后直接将两者相加。 具体运算如下式展示, 其中 \({F} \in  {\mathbb{R}}^{{B} * {N} * \text{ out }}\) 为时空注意力层、时空卷积层的输出, out 为输出通道数(根据上文为 64), \({I} \in  {\mathbb{R}}^{{B} * {N} * {{in}}}\) 为外部信息, \({{in}}\) 为输入外部信息通道数(SZ 为 2, 融合天气和 POI 信息, LA 为 1, 融合道路占有率指标)。 \({W} \in  {\mathbb{R}}^{{{in}} * {{out}}}\)。
\begin{equation}
{H} = {F} + {I} * {W}.
\end{equation}

另外通过引入外源信息, 能够有效地增强模型对时间和空间注意力矩阵的学习能力。多信息融合反映了更多节点的信息, 有利于增强模型对节点的相似性学习。

SZ-taxi 数据集在后续的研究中收集了额外的补充了道路节点的 POI 属性和 31 天中的天气状况, 得到本研究最终的数据集。图 \ref{fig:weather-speed-comparison} 展示了晴天和雨天, 随机选择一些节点的行车速度均值, 可以发现晴天和雨天的速度会有明显差异。

\begin{figure}[htbp]
  \centering
  \includegraphics[width=0.6\textwidth]{01960c08-a17d-7162-9acb-7f9eeb4038d6_30_479_305_742_553_0.jpg}
  \caption{不同天气状况下速度对比}
  \label{fig:weather-speed-comparison}
\end{figure}

实验结果如表 \ref{tab:multi} 所示, 该结果可以论证, 时空注意力机制-图网络也可以拓展到多信息融合, 另外融合多信息可以在一定程度上提高预测精度, LA 数据集在融合进道路占有率指标后（道路上被占用的车辆或其 他交通参与者数量与道路总容量的比值）, RMSE 下降 3.47\%。因此, 什么样的信息将有助于流量预测, 如何更加有效利用外源信息, 并将其更好融合进模型, 也是未来研究的方向。


\begin{table}
  \caption{融合多信息与原始数据预测结果}
  \label{tab:multi}
  \centering
  \begin{tblr}{
    hline{1, Z} = {1pt}, 
    hline{2}, 
    width = \textwidth, 
    colspec = {X[2, c]X[1, c]X[1, c]X[1, c]}, 
  }
    & RMSE & MAE & MAPE \\
    SZ with Weather \& POI & 4.02 & 2.71 & 9.21\% \\
    SZ & 4.05 & 2.77 & 9.21\% \\
    LA with Occupancy Rate & 6.92 & 3.64 & 10.85\% \\
    LA & 7.17 & 3.67 & 11.84\% \\
  \end{tblr}
\end{table}
